\subsection*{Tileflow Program: Orchestrating the Data Movement}

The \emph{Tileflow Program} consists of \emph{Croqtile functions}. As described earlier, it manages the movement of data between the host and the target device, ensuring that data is copied correctly across different storage locations.

For convenience, let's revisit the code:

\begin{lstlisting}[language=croqtile,caption={The tileflow program}]
__co__ s32 [6, 17, 128] ele_add(s32 [6, 17, 128] lhs, s32 [6, 17, 128] rhs) {
  s32 [lhs.span] output; // Use same shape as lhs
  parallel p by 6 {
    with index in [17, 4] { // Tiling factors
      foreach index {
        lhs_load = dma.copy lhs.chunkat(p, index) => local;
        rhs_load = dma.copy rhs.chunkat(p, index) => local;
        local s32 [lhs_load.span] l1_out;
        // Call kernel with loaded data
        call kernel(lhs_load.data, rhs_load.data, l1_out, |lhs_load.span|);
        // Store result back to output
        dma.copy l1_out => output.chunkat(p, index);
      }
    }
  }
  return output;
}
\end{lstlisting}

The \texttt{\_\_co\_\_} function accepts \texttt{lhs} and \texttt{rhs}, both with shape \texttt{[6, 17, 128]} and element type \texttt{s32} (signed 32-bit integer). The output has the same shape and element type.

The \texttt{parallel p by 6 \{.\allowbreak{}.\allowbreak{}.\allowbreak{}\}} block indicates that the enclosed code runs in parallel. Specifically, six instances of the code are executed concurrently. This implies a transition in the execution environment from the host to the device.

Inside the \texttt{parallel-by} block, \texttt{with-in} binds \texttt{index} to the bounds \texttt{[17, 4]}. This two-dimensional bounded tuple supplies the indices for the \texttt{foreach} loop.

The \texttt{foreach index \{.\allowbreak{}.\allowbreak{}.\allowbreak{}\}} statement is equivalent to the following C code:

\begin{lstlisting}[language={},caption={The equivalent C++ loop}]
for (int x = 0; x < 17; x++)  // assume 'x' represents the 1st element of 'index'
  for (int y = 0; y < 4; y++) { ... }  // and 'y' represents the 2nd element of 'index'
\end{lstlisting}

Within the \texttt{foreach} block, the \texttt{dma.copy} statement describes how data movement occurs. For example, consider the statement \texttt{lhs\_\allowbreak{}load = dma.\allowbreak{}copy lhs.\allowbreak{}chunkat(p, index) =\textgreater{} local;}:

\begin{itemize}[leftmargin=1.6em]
\item The symbol \texttt{lhs\_load} is called \textbf{future} of the DMA operation, which contains information about the DMA destination.
\item \texttt{dma.copy} invokes a direct DMA data transfer without transforming the shape of the data. The expression on the left-hand side of \texttt{=\textgreater{}} represents the DMA source, while the right-hand side represents the destination.
\item In this case, the destination is specified as a \texttt{local} buffer, which is automatically allocated by the Croqtile compiler.
\item The source expression \texttt{lhs.\allowbreak{}chunkat(p, index)} selects a chunk of \texttt{lhs}. The bounds of \texttt{p} and \texttt{index}, namely \texttt{[6, 17, 4]}, determine the partition. Dividing the input shape \texttt{[6, 17, 128]} by these bounds gives chunk shape \texttt{[1, 1, 32]}. The current values select the chunk: instance \texttt{p = 1} and iteration \texttt{index = \{16, 2\}} select chunk \texttt{(1, 16, 2)}, whose element origin is \texttt{(1, 16, 64)}.
\end{itemize}

The following figure shows how the bounds determine the chunk shape and the indices select a chunk.

Figure~\ref{fig:f14-chunkat_3d} shows how the bounds determine the chunk shape.

\begin{figure}[htbp]\centering
\includegraphics[width=\linewidth,height=0.56\textheight,keepaspectratio]{f14-chunkat_3d_light.png}
\caption{Dividing shape [6, 17, 128] by bounds [6, 17, 4] gives chunks of shape [1, 1, 32]. Chunk (1, 16, 2) starts at element (1, 16, 64).}
\label{fig:f14-chunkat_3d}
\end{figure}

With the DMA statement, different chunks of data tiled from \texttt{lhs} are moved iteratively and in parallel from the host to the device's \texttt{local} storage. Similarly, the DMA statement manages \texttt{rhs} by moving it in small chunks to \texttt{local} storage for processing.

The declaration \texttt{local s32 [lhs\_\allowbreak{}load.\allowbreak{}span] l1\_\allowbreak{}out;} creates a buffer for each parallel instance, with the same shape as the loaded chunk. The subsequent \texttt{call} invokes the device function \texttt{kernel}, which writes its result to that buffer. A final DMA statement copies the result back to the host output, completing the iteration.

In this code, each parallel thread runs \texttt{17x4} iterations, with each iteration handling a \texttt{1x1x32}-sized chunk of data. The \emph{croqtile program} terminates when all \texttt{6} parallel threads have completed their iterations. It then returns the output buffer to its caller, the host program.
